% GENERATED FILE. Edit canonical lessons, then run npm run generate:books.
% canonical-source-hash: fnv1a64:2f0904333d96a062
% canonical-lessons: KA-C261-sakuprani, KA-C261-karu, KA-C261-alilu, KA-C261-halli, KA-C261-jirale, KA-R261-first-pass-words-from-chapters-253-256, KA-R261-second-pass-words-from-chapters-257-261

\chapter{Pet, Calf, Squirrel, Lizard, Cockroach}
\label{ch:ka-pet-calf-squirrel-lizard-cockroach}
\hlchaptermodality{261}

\begin{chapteropening}
\textbf{By the end of this chapter:} \emph{I can say a pet, a calf, a squirrel, a lizard and a cockroach.}

Complete the last lesson of chapter 261: I can say a pet, a calf, a squirrel, a lizard and a cockroach.
\end{chapteropening}

\section[\texorpdfstring{\kn{ಸಾಕುಪ್ರಾಣಿ} (sākuprāṇi)}{ಸಾಕುಪ್ರಾಣಿ (sākuprāṇi)}]{\kn{ಸಾಕುಪ್ರಾಣಿ} (sākuprāṇi) — a pet}

\label{lesson:KA-C261-sakuprani}



\begin{quote}
Before the new one: say the Kannada for a vehicle, then the Kannada for an autorickshaw.
\end{quote}

\subsection*{You'll want to know: \kn{ಸಾಕುಪ್ರಾಣಿ}}
\textbf{\kn{ಸಾಕುಪ್ರಾಣಿ}} — \emph{sākuprāṇi} — \textquotedblleft{}a pet\textquotedblright{}.

\begin{grammarlens}[title={What you've built}]
One more word for this chapter.
\end{grammarlens}

\subsection*{Your turn}
\begin{itemize}
\raggedright
  \item \emph{Say it:} \emph{sākuprāṇi}
  \item \emph{Say it:} \emph{sākuprāṇi}, once more
  \item \emph{Say it:} say \emph{sākuprāṇi}
  \item \emph{Recall:} say the Kannada for a vehicle, then the Kannada for an autorickshaw, then say \emph{sākuprāṇi} again
\end{itemize}

\begin{tcolorbox}[breakable,colback=teal!4,colframe=teal!35!black,title={Before you move on}]
What does \kn{ಸಾಕುಪ್ರಾಣಿ} mean? (\textquotedblleft{}A pet\textquotedblright{}.) Say it once more.
\end{tcolorbox}
\section[\texorpdfstring{\kn{ಕರು} (karu)}{ಕರು (karu)}]{\kn{ಕರು} (karu) — a calf}

\label{lesson:KA-C261-karu}



\begin{quote}
Before the new one: say the Kannada for an autorickshaw, then the Kannada for a pet.
\end{quote}

\subsection*{You'll want to know: \kn{ಕರು}}
\textbf{\kn{ಕರು}} — \emph{karu} — \textquotedblleft{}a calf\textquotedblright{}.

\begin{grammarlens}[title={What you've built}]
One more word for this chapter.
\end{grammarlens}

\subsection*{Your turn}
\begin{itemize}
\raggedright
  \item \emph{Say it:} \emph{karu}
  \item \emph{Say it:} \emph{karu}, once more
  \item \emph{Say it:} say \emph{karu}
  \item \emph{Recall:} say the Kannada for an autorickshaw, then the Kannada for a pet, then say \emph{karu} again
\end{itemize}

\begin{tcolorbox}[breakable,colback=teal!4,colframe=teal!35!black,title={Before you move on}]
What does \kn{ಕರು} mean? (\textquotedblleft{}A calf\textquotedblright{}.) Say it once more.
\end{tcolorbox}
\section[\texorpdfstring{\kn{ಅಳಿಲು} (aḷilu)}{ಅಳಿಲು (aḷilu)}]{\kn{ಅಳಿಲು} (aḷilu) — a squirrel}

\label{lesson:KA-C261-alilu}



\begin{quote}
Before the new one: say the Kannada for a pet, then the Kannada for a calf.
\end{quote}

\subsection*{You'll want to know: \kn{ಅಳಿಲು}}
\textbf{\kn{ಅಳಿಲು}} — \emph{aḷilu} — \textquotedblleft{}a squirrel\textquotedblright{}.

\begin{grammarlens}[title={What you've built}]
One more word for this chapter.
\end{grammarlens}

\subsection*{Your turn}
\begin{itemize}
\raggedright
  \item \emph{Say it:} \emph{aḷilu}
  \item \emph{Say it:} \emph{aḷilu}, once more
  \item \emph{Say it:} say \emph{aḷilu}
  \item \emph{Recall:} say the Kannada for a pet, then the Kannada for a calf, then say \emph{aḷilu} again
\end{itemize}

\begin{tcolorbox}[breakable,colback=teal!4,colframe=teal!35!black,title={Before you move on}]
What does \kn{ಅಳಿಲು} mean? (\textquotedblleft{}A squirrel\textquotedblright{}.) Say it once more.
\end{tcolorbox}
\section[\texorpdfstring{\kn{ಹಲ್ಲಿ} (halli)}{ಹಲ್ಲಿ (halli)}]{\kn{ಹಲ್ಲಿ} (halli) — a lizard}

\label{lesson:KA-C261-halli}



\begin{quote}
Before the new one: say the Kannada for a calf, then the Kannada for a squirrel.
\end{quote}

\subsection*{You'll want to know: \kn{ಹಲ್ಲಿ}}
\textbf{\kn{ಹಲ್ಲಿ}} — \emph{halli} — \textquotedblleft{}a lizard\textquotedblright{}.

\begin{grammarlens}[title={What you've built}]
One more word for this chapter.
\end{grammarlens}

\subsection*{Your turn}
\begin{itemize}
\raggedright
  \item \emph{Say it:} \emph{halli}
  \item \emph{Say it:} \emph{halli}, once more
  \item \emph{Say it:} say \emph{halli}
  \item \emph{Recall:} say the Kannada for a calf, then the Kannada for a squirrel, then say \emph{halli} again
\end{itemize}

\begin{tcolorbox}[breakable,colback=teal!4,colframe=teal!35!black,title={Before you move on}]
What does \kn{ಹಲ್ಲಿ} mean? (\textquotedblleft{}A lizard\textquotedblright{}.) Say it once more.
\end{tcolorbox}
\section[\texorpdfstring{\kn{ಜಿರಳೆ} (jiraḷe)}{ಜಿರಳೆ (jiraḷe)}]{\kn{ಜಿರಳೆ} (jiraḷe) — a cockroach}

\label{lesson:KA-C261-jirale}



\begin{quote}
Before the new one: say the Kannada for a squirrel, then the Kannada for a lizard.
\end{quote}

\subsection*{You'll want to know: \kn{ಜಿರಳೆ}}
\textbf{\kn{ಜಿರಳೆ}} — \emph{jiraḷe} — \textquotedblleft{}a cockroach\textquotedblright{}.

\begin{grammarlens}[title={What you've built}]
One more word for this chapter.
\end{grammarlens}

\subsection*{Your turn}
\begin{itemize}
\raggedright
  \item \emph{Say it:} \emph{jiraḷe}
  \item \emph{Say it:} \emph{jiraḷe}, once more
  \item \emph{Say it:} say \emph{jiraḷe}
  \item \emph{Recall:} say the Kannada for a squirrel, then the Kannada for a lizard, then say \emph{jiraḷe} again
\end{itemize}

\begin{tcolorbox}[breakable,colback=teal!4,colframe=teal!35!black,title={Before you move on}]
What does \kn{ಜಿರಳೆ} mean? (\textquotedblleft{}A cockroach\textquotedblright{}.) Say it once more.
\end{tcolorbox}
\section[(dialogue)]{Words from chapters 253-256 — a first pass}

\label{lesson:KA-R261-first-pass-words-from-chapters-253-256}



\begin{quote}
Say the words for a lizard and a cockroach.
\end{quote}

\subsection*{Your turn}
Say these aloud:

\begin{itemize}
\raggedright
  \item an example, an event, a memory, life, a rule
  \item thin, unwell, sharp, brown, pink
  \item naughty, calm, selfish, serious, round
  \item hollow, flat, crooked, blurred, crowded
\end{itemize}

\begin{tcolorbox}[breakable,colback=teal!4,colframe=teal!35!black,title={Before you move on}]
An example? (\textbf{\kn{ಉದಾಹರಣೆ}}, \emph{udāharaṇe}.) An event? (\textbf{\kn{ಘಟನೆ}}, \emph{ghaṭane}.) A memory? (\textbf{\kn{ನೆನಪು}}, \emph{nenapu}.) Life? (\textbf{\kn{ಜೀವನ}}, \emph{jīvana}.) A rule? (\textbf{\kn{ನಿಯಮ}}, \emph{niyama}.) Thin? (\textbf{\kn{ತೆಳು}}, \emph{teḷu}.) Unwell? (\textbf{\kn{ಅಸ್ವಸ್ಥ}}, \emph{asvastha}.) Sharp? (\textbf{\kn{ಚೂಪಾದ}}, \emph{cūpāda}.) Brown? (\textbf{\kn{ಕಂದು}}, \emph{kandu}.) Pink? (\textbf{\kn{ಗುಲಾಬಿ}}, \emph{gulābi}.) Naughty? (\textbf{\kn{ತುಂಟ}}, \emph{tuṇṭa}.) Calm? (\textbf{\kn{ಶಾಂತ}}, \emph{śānta}.) Selfish? (\textbf{\kn{ಸ್ವಾರ್ಥಿ}}, \emph{svārthi}.) Serious? (\textbf{\kn{ಗಂಭೀರ}}, \emph{gambhīra}.) Round? (\textbf{\kn{ಗುಂಡಗೆ}}, \emph{guṇḍage}.) Hollow? (\textbf{\kn{ಟೊಳ್ಳು}}, \emph{ṭoḷḷu}.) Flat? (\textbf{\kn{ಚಪ್ಪಟೆ}}, \emph{cappaṭe}.) Crooked? (\textbf{\kn{ಡೊಂಕು}}, \emph{ḍoṅku}.) Blurred? (\textbf{\kn{ಮಸುಕು}}, \emph{masuku}.) Crowded? (\textbf{\kn{ಕಿಕ್ಕಿರಿದ}}, \emph{kikkirida}.)
\end{tcolorbox}
\section[(dialogue)]{Words from chapters 257-261 — a second pass}

\label{lesson:KA-R261-second-pass-words-from-chapters-257-261}



\begin{quote}
Say the words for a tap and stairs.
\end{quote}

\subsection*{Your turn}
Say these aloud:

\begin{itemize}
\raggedright
  \item a tap, stairs, flour, lentils, tamarind
  \item buttermilk, butter, sweet milk pudding, a biscuit, bread
  \item a button, a sweater, a shawl, a handkerchief, a necklace
  \item luggage, a key, a newspaper, a vehicle, an autorickshaw
  \item a pet, a calf, a squirrel, a lizard, a cockroach
\end{itemize}

\begin{tcolorbox}[breakable,colback=teal!4,colframe=teal!35!black,title={Before you move on}]
A tap? (\textbf{\kn{ನಲ್ಲಿ}}, \emph{nalli}.) Stairs? (\textbf{\kn{ಮೆಟ್ಟಿಲು}}, \emph{meṭṭilu}.) Flour? (\textbf{\kn{ಹಿಟ್ಟು}}, \emph{hiṭṭu}.) Lentils? (\textbf{\kn{ಬೇಳೆ}}, \emph{bēḷe}.) Tamarind? (\textbf{\kn{ಹುಣಸೆಹಣ್ಣು}}, \emph{huṇasehaṇṇu}.) Buttermilk? (\textbf{\kn{ಮಜ್ಜಿಗೆ}}, \emph{majjige}.) Butter? (\textbf{\kn{ಬೆಣ್ಣೆ}}, \emph{beṇṇe}.) Sweet milk pudding? (\textbf{\kn{ಪಾಯಸ}}, \emph{pāyasa}.) A biscuit? (\textbf{\kn{ಬಿಸ್ಕತ್ತು}}, \emph{biskattu}.) Bread? (\textbf{\kn{ಬ್ರೆಡ್}}, \emph{breḍ}.) A button? (\textbf{\kn{ಗುಂಡಿ}}, \emph{guṇḍi}.) A sweater? (\textbf{\kn{ಸ್ವೆಟರ್}}, \emph{sveṭar}.) A shawl? (\textbf{\kn{ಶಾಲು}}, \emph{śālu}.) A handkerchief? (\textbf{\kn{ಕರವಸ್ತ್ರ}}, \emph{karavastra}.) A necklace? (\textbf{\kn{ಸರ}}, \emph{sara}.) Luggage? (\textbf{\kn{ಸಾಮಾನು}}, \emph{sāmānu}.) A key? (\textbf{\kn{ಕೀಲಿ}}, \emph{kīli}.) A newspaper? (\textbf{\kn{ಪತ್ರಿಕೆ}}, \emph{patrike}.) A vehicle? (\textbf{\kn{ವಾಹನ}}, \emph{vāhana}.) An autorickshaw? (\textbf{\kn{ಆಟೋ}}, \emph{āṭō}.) A pet? (\textbf{\kn{ಸಾಕುಪ್ರಾಣಿ}}, \emph{sākuprāṇi}.) A calf? (\textbf{\kn{ಕರು}}, \emph{karu}.) A squirrel? (\textbf{\kn{ಅಳಿಲು}}, \emph{aḷilu}.) A lizard? (\textbf{\kn{ಹಲ್ಲಿ}}, \emph{halli}.) A cockroach? (\textbf{\kn{ಜಿರಳೆ}}, \emph{jiraḷe}.)
\end{tcolorbox}
